\chapter{\textbf{Introduction}}
\label{chap:introduction}

\section*{Outline}
\addcontentsline{toc}{section}{\texorpdfstring{\hskip14mm}{} Outline}

This chapter tells you what \texttt{Cast3M} is, how to install it, and how to
run your first file. If you only want to get a \texttt{.dgibi} file supplied
with a lecture to run on your own machine, read
sections~\ref{sec:intro:install} and~\ref{sec:intro:running} and skip the rest
until you need it.

\section{Historical remarks}

The CEA -- Commissariat à l'énergie atomique et aux énergies alternatives,
Saclay, France (\url{https://www.cea.fr}) is a French government-funded
technological research organisation, which has developed over the past forty
years a dedicated toolbox for finite element computations: \texttt{Cast3M}.
The toolbox contains the essential ingredients needed to perform a finite
element computation. Its main application area is mechanics -- elasticity,
elasto-visco-plasticity, and analyses adapted to equilibrium, buckling,
vibration or dynamics problems. Other application areas include thermal,
hydraulic and electromagnetic analysis. Its modular character makes it easy to
program new applications, which makes the tool well adapted to coupled
multiphysics problems.

\texttt{Cast3M} is built on a family of \emph{objects}, i.e. data structures,
and a lexicon of \emph{operators}, i.e. elementary operations acting on those
objects. Together they form a specific programming language called
\textsc{gibiane}. The code itself is written in an extended Fortran called
\textsc{esope} which permits direct handling of the specific data structures.

\section{The Cast3M website}
\label{sec:intro:website}

Essentially everything you will need is on
\url{https://www-cast3m.cea.fr/}. It is worth spending ten minutes finding
your way around it before you start, because the site -- not this booklet --
is the authoritative and up-to-date reference. The pages you will use most
are:

\begin{center}
\fcolorbox{black}{Gray}{%
\begin{tabular}{@{\hspace{2mm}}l@{\hspace{4mm}}p{72mm}@{\hspace{2mm}}}
\multicolumn{2}{c}{\textbf{Where to look on the Cast3M site}} \\[2.mm]
\textbf{Notices} & the reference manual: one page per operator, about 1\,100 of
them, with a full-text search over the notices \\[1.mm]
\texttt{?page=notices} & \\[2.mm]
\textbf{Exemples} & roughly 1\,600 searchable example files -- usually the
fastest way to learn an operator \\[1.mm]
\texttt{?page=exemples} & \\[2.mm]
\textbf{Procédures} & the $\approx$560 standard procedures written in
\textsc{gibiane} \\[1.mm]
\texttt{?page=procedures} & \\[2.mm]
\textbf{Téléchargement} & download and licence form \\[1.mm]
\texttt{?xml=download1} & \\[2.mm]
\textbf{Formations} & free CEA training sessions, and the course slides \\[1.mm]
\texttt{?xml=formations} & \\[2.mm]
\textbf{Utilitaires} & editor syntax files and the mesh conversion scripts \\[1.mm]
\texttt{?xml=utilitaires} & \\[2.mm]
\textbf{FAQ} & \texttt{?page=faqphp} \\
\end{tabular}}
\end{center}

\noindent prefix each of these with
\texttt{https://www-cast3m.cea.fr/index.php}, for instance
\url{https://www-cast3m.cea.fr/index.php?page=notices}. A single operator
notice can be reached directly, which is handy to bookmark:
\url{https://www-cast3m.cea.fr/index.php?page=notices&notice=LIRE}.

Two further resources did not exist when the first version of these notes was
written and are worth knowing about:

\begin{itemize}
\item the \textbf{theory documentation}, rewritten as a modern searchable site
  at \url{https://doc-cast3m.readthedocs.io/} (in French), covering
  quasi-static and dynamic mechanics, concrete models, transient thermal
  analysis and topology optimisation;
\item the \textbf{user mailing list}, which is where questions actually get
  answered. Write to
  \texttt{cast3m-util@\allowbreak umontpellier.fr}; to subscribe, send
  a message to \texttt{sympa@\allowbreak umontpellier.fr} containing the single line
  \texttt{SUB cast3m-util \textit{Surname} \textit{Firstname}}, or use the web
  interface \url{https://listes.umontpellier.fr/sympa}. For problems with the
  code itself there is \texttt{support-cast3m@cea.fr}.
\end{itemize}

\noindent Note that the operator notices are written in French. Several
documents aimed at beginners are available in English, in particular
\emph{Starting with Cast3M} and \emph{The PASAPAS procedure}, both linked from
the \textbf{Formations} page; an older but still useful introduction is
Fichoux~\cite{Fichoux:1999}.

\section{Obtaining and installing Cast3M}
\label{sec:intro:install}

\subsection*{Version numbering -- read this first}

Releases are named after the year: the current one is \textbf{Cast3M 2026},
internal version number \texttt{26.0}. Two consequences matter in daily use:

\begin{itemize}
\item \textbf{the executable carries the two-digit year}: you launch
  \texttt{castem26}, not \texttt{castem}. A script written for an older
  release will call \texttt{castem22} or \texttt{castem19} and will simply not
  be found on your machine;
\item \textbf{the environment variables carry it too}:
  \texttt{CASTEM\_PROCEDUR26}, \texttt{CASTEM\_NOTICE26}.
\end{itemize}

\noindent In practice one release appears per year. Throughout this booklet
\texttt{castem26} stands for whichever version you have installed -- substitute
your own number.

\subsection*{Licence and download}

\texttt{Cast3M} is \textbf{free of charge for teaching and public research},
but it is \emph{not} open-source software and it is not anonymous to download.
You fill in a short form (name, email, institution, status -- student,
doctoral candidate, teacher, \ldots\ -- and intended use) at
\url{https://www-cast3m.cea.fr/index.php?xml=download1} and accept the
\emph{Contrat de licence Cast3M « Éducation \& Recherche »}. The practical
points of that licence:

\begin{itemize}
\item use is restricted to teaching and research; providing services to third
  parties, and any commercial or advertising use, is excluded (a commercial
  licence exists separately);
\item the licence runs for five years and covers one workstation, or one
  server accessible to a single research or teaching team;
\item you may not pass the software on to anyone else -- \textbf{send your
  students to the download page, do not redistribute the archive};
\item the sources are supplied and you may modify them, but anything you
  submit back for integration is assigned to the CEA;
\item publications using \texttt{Cast3M} must cite it and name the CEA as the
  holder of the intellectual property rights.
\end{itemize}

\subsection*{Installation}

The distribution is a self-contained archive with an installer script inside;
there is nothing to compile. Unpack it, read the \texttt{Readme.txt} /
\texttt{Lisez-moi.txt} it contains, and run the installer:

\begin{center}
\fcolorbox{black}{Gray}{%
\begin{tabular}{@{\hspace{2mm}}l@{\hspace{4mm}}l@{\hspace{4mm}}p{58mm}@{\hspace{2mm}}}
\multicolumn{3}{c}{\textbf{Installation}} \\[2.mm]
\textbf{Platform} & \textbf{Installer} & \textbf{Notes} \\[1.mm]
GNU/Linux & \texttt{INSTALL\_CAST3M.sh} & \texttt{tar} archive; run the script \\[1.mm]
Windows   & \texttt{INSTALL\_CAST3M.bat} & \texttt{zip} archive; double-click \\[1.mm]
macOS     & \texttt{INSTALL\_CAST3M.sh} & \texttt{tar} archive; needs Xcode, then
                                          \texttt{sudo xcodebuild -license} \\
\end{tabular}}
\end{center}

\noindent A graphical installer then takes over and sets the environment up
for you. Three remarks:

\begin{itemize}
\item \textbf{macOS is now officially supported.} On Apple Silicon (ARM64)
  the current version is available like on the other platforms -- no virtual
  machine, no Docker, no self-compilation. The situation described in earlier
  editions of these notes, where Mac users relied on unofficial builds, is
  obsolete. Intel Macs, however, are frozen at version \texttt{22.0}.
\item \textbf{There is no official Docker image and no conda package}, despite
  what you may find in forum posts.
\item On Windows, older notes told you to create \texttt{C:\textbackslash tmp}
  by hand and to set a \texttt{CASTEM} variable. The current installer handles
  this; only fall back on manual settings if it fails.
\end{itemize}

\subsection*{Environment variables}

You should not normally have to set these yourself, but they are what the
installer configures and what you will need if you move an installation or run
on a cluster:

\begin{center}
\fcolorbox{black}{Gray}{%
\begin{tabular}{@{\hspace{2mm}}>{\ttfamily}l@{\hspace{4mm}}p{95mm}@{\hspace{2mm}}}
\multicolumn{2}{c}{\textbf{Environment variables}} \\[2.mm]
CASTEM\_PROCEDUR26 & directory the standard procedures are read from \\
CASTEM\_NOTICE26   & directory the operator notices are read from \\
ESOPE\_PARAM       & size of the \textsc{esope} central memory zone, in machine
                     words \\
\end{tabular}}
\end{center}

\noindent Procedures are looked up in the current directory first, then in
\texttt{./procedur}, then in the installation directory -- which is why you can
override a standard procedure simply by dropping your own copy next to your
\texttt{.dgibi} file.

If a long computation stops with a message of the form

\begin{center}\texttt{--- GEMAT ERROR --- \ldots\ SEGINI, IBLIS --- PAS ASSEZ DE
PLACE EN MEMOIRE}\end{center}

\noindent it has run out of \textsc{esope} memory. Increase
\texttt{ESOPE\_PARAM}, but keep it below your physical RAM or the machine will
swap and the run will become slower than useless. Inside a file the related
settings are \op{opti 'PLAC' n;} (minimum free memory to maintain, default
$P_{tot}/10$) and \op{opti 'NGMA' n;} (maximum matrix size in words, default
$100\,000$).

\medskip
\noindent\emph{A note for readers of the older edition:} the
\texttt{castem -m XL} memory option described there no longer appears in the
current documentation. Use \texttt{ESOPE\_PARAM} instead.

\section{Running Cast3M}
\label{sec:intro:running}

There are two ways to run the program, and you will use both.

\begin{castemcom}
Launching \texttt{castem26} with a file name executes the commands in that
file and then \textbf{returns the hand to you} in an interactive
\textsc{gibiane} session, where you can go on typing commands against the
objects the file has just built. This is the normal way to work: put the
computation in a file, then explore the results interactively.

Technically the program copies \texttt{myfile.dgibi} into \texttt{fort.3}
(logical unit 3) and reads it, then reads what you type from \texttt{fort.98}
(logical unit 98).

\medskip
Two consequences worth remembering. First, after a crash \texttt{fort.98}
still holds everything you typed, which is often the quickest way to recover
your work -- but it is wiped at the next start. Second, both files live in the
working directory, so \textbf{do not run two Cast3M sessions in the same
directory} at the same time: they will fight over the same files.
\end{castemcom}
\begin{castemlines}
\begin{verbatim}
$ castem26 myfile.dgibi
\end{verbatim}
\end{castemlines}

\begin{castemcom}
Launched without an argument, the program goes straight to the interactive
prompt. Each line you type is echoed back prefixed with \texttt{*}.

\medskip
Useful while you are getting started, and for trying out a single operator,
but anything you want to keep belongs in a file.
\end{castemcom}
\begin{castemlines}
\begin{verbatim}
$ castem26
 ...
$ a = 3;
* a = 3;
$
\end{verbatim}
\end{castemlines}

\noindent Once the program is running, a few commands are worth memorising:

\begin{itemize}
\item \op{fin;} stops the program;
\item \op{opti donn 5;} placed inside \texttt{myfile.dgibi} interrupts the
  reading of the file at that point and gives you the hand -- the single most
  useful debugging trick in \textsc{gibiane}, since it lets you stop in the
  middle of a computation and inspect the objects built so far;
\item \op{opti donn 3;} typed at the prompt hands control back to the file;
\item \op{info \textit{operator};} prints the notice of an operator in the
  terminal (in French). \op{aide} and \op{noti} do related things.
\end{itemize}

\noindent A further useful invocation is \texttt{castem26 -test}, which runs
the test-case suite -- a quick way to confirm that an installation is sane.

\section{Editing gibiane files}
\label{sec:intro:editing}

A \texttt{.dgibi} file is plain ASCII text, so any editor will do, but syntax
highlighting catches a surprising number of mistakes before you run anything.
Configuration files for a long list of editors are on the
\textbf{Utilitaires} page of the website. Two are worth singling out:

\begin{itemize}
\item \textbf{Visual Studio Code}, with the \texttt{gibiane-vscode} extension
  (\url{https://github.com/charles-zablit/gibiane-vscode}), which adds
  highlighting, completion and go-to-definition. This is the recommendation
  for anyone starting today;
\item \textbf{Vim}, \textbf{Emacs}, \textbf{Notepad++}, \textbf{Kate},
  \textbf{gedit} and \textbf{Geany} all have published syntax files.
\end{itemize}

\noindent Whatever you use, \textbf{switch off tab characters}: a literal
\texttt{TAB} in a \texttt{.dgibi} file causes reading errors that are tedious
to diagnose.

\section{Structure of the program}

From the programming point of view \texttt{Cast3M} has two levels:

\begin{itemize}
\item a compiled low-level language, \textsc{esope}. This is invisible to the
  normal user; it is the language of the authors of \texttt{Cast3M}. It is a
  derivative of \textsc{fortran} enriched with operations for the manipulation
  of specific objects (creation, copy, removal, \ldots\ of data structures).
\item a high-level interpreted language, \textsc{gibiane}, formed by a lexicon
  of operations, which is the language of the normal user and the one
  exemplified throughout this text.
\end{itemize}

\noindent From this point of view \texttt{Cast3M} is comparable with
\textsc{Matlab}, \textsc{Scilab}, \textsc{Octave}, \textsc{Mathematica} or
\textsc{Maple}: \textsc{esope} is the language of the creator and
\textsc{gibiane} the language of the user.

Two extension routes exist if \textsc{gibiane} is not enough. New operators
are written in \textsc{esope}. New \emph{material laws}, on the other hand,
are nowadays written with \textbf{MFront}
(\url{https://thelfer.github.io/tfel/web}), an open-source code generator
developed at the CEA and bundled with \texttt{Cast3M}; it plugs in through the
generic \texttt{UMAT} interface and saves you from writing \textsc{esope} for
what is the most common reason to want to extend a finite element code.

The essence of \textsc{gibiane} can be summarised in one line:

\begin{castemcom}
The \texttt{;} marks the end of the command and makes the system execute the
operator with \texttt{objet\_1} as its argument, creating a new object
\texttt{objet\_2} as the result.
\end{castemcom}
\begin{castemlines}
\begin{verbatim}
objet_2 = operator objet_1 ;
\end{verbatim}
\end{castemlines}

\noindent The basic syntax rules of \textsc{gibiane} are:

\begin{itemize}
\item operator names are defined by their first 4 characters, so \op{vibr}
  denotes the \op{vibr[ation]} operator; spelling the name out in full is
  still permitted;
\item object names should be at most 8 characters, and their first 4
  characters should not coincide with an operator name -- that would shadow
  the operator;
\item a statement may run over several lines, with no continuation character
  needed, up to a limit of 500 characters in total. Individual lines should
  stay within 72 characters;
\item \textbf{there is no operator precedence}: operations are executed in
  order of appearance, so
  \[ \op{a + b * c} \;=\; \op{(a + b) * c} \;\neq\; \op{a + (b * c)} \]
  Use brackets systematically in algebraic expressions. This is the single
  most common source of silently wrong results for newcomers;
\item upper and lower case are not distinguished, except inside quoted
  strings: \op{vibr}, \op{Vibr} and \op{VIBR} are the same operator, but
  \texttt{'SMYY'} and \texttt{'smyy'} are not the same string;
\item a \texttt{*} in the first column starts a comment;
\item tab characters and invisible characters inserted by some editors must be
  avoided, as they create file-reading errors.
  \sindex{editors (editeurs, outils)}
\end{itemize}

\subsection*{Objects}

Objects are numerical or logical values organised in data structures
describing mathematical and physical fields. The classical ones describing
numbers are:

\begin{itemize}
\item \op{entier} \dindex{entier}{integer} is an integer, written without a
  \texttt{.};
\item \op{flottant} \dindex{flottant}{real number} is a real number, written
  either as \texttt{123.4} or as \texttt{1.234e2} $=1.234\times10^{2}$;
\item \op{logique} \dindex{logique}{logical value} are logical objects used in
  boolean operations, with the two values \texttt{VRAI} (true) and
  \texttt{FAUX} (false).
\end{itemize}

\noindent Numbers and strings can be organised in lists:

\begin{itemize}
\item \op{listenti[er]} \dindex{listentier}{list of integers} a list of
  integers, assembled with \op{lect[ure]}: \sindex{lect[ure]}
  \op{mylist = lect 1 6 89;}
\item \op{listreel} \dindex{listreel}{list of real numbers} a list of reals,
  assembled with \op{prog[ramer]}: \sindex{prog[ramer]}
  \op{mylist = prog 1. 6.e3 8.9;}
\item \op{listmots} \dindex{listmots}{list of strings} a list of strings
  (fr. \emph{mots}, words), assembled with \op{mots}: \sindex{mots}
  \op{mylist = mots 'TITI' 'Toto' 'tete';}
\end{itemize}

\begin{optable}{Creating lists of objects}
lect & lists of integers \\
prog & lists of reals \\
mots & lists of strings \\
table & generalised list, indexed by integers or strings \\
\end{optable}

\noindent The \op{table} deserves a special mention because it is used
everywhere in practice: it is a generalised list whose indices may be integers
or strings, and it is how \op{pasapas} returns the results of a nonlinear
computation (chapter~\ref{chap:plasticity}).

Another class of objects is specific to a finite element computation:

\begin{itemize}
\item the \op{maillage} \dindex{maillage}{mesh} (mesh) is described by an
  element type -- \texttt{SEG2}, \texttt{TRI3}, \texttt{QUA4}, \ldots\ -- a
  list of nodes and their connectivity;
\item a \op{chpoint} \dindex{chpoint}{nodal field}, i.e. \emph{champ par
  point}, is a scalar, vector or tensor field defined by its values at the
  nodes of a mesh: temperature, displacement, \ldots;
\item a \op{mchaml} \dindex{mchaml}{field on elements}, i.e. \emph{champ par
  élément}, is a field defined by its values at specific points of each
  element -- typically the Gauss integration points, where stresses, strains
  and temperature gradients live;
\item a \op{rigidite} \dindex{rigidite}{stiffness} (stiffness) is a stiffness
  matrix, possibly carrying a set of preconditioners.
\end{itemize}

\noindent The distinction between \op{chpoint} and \op{mchaml} is the one
beginners stumble over most often: a displacement is a \op{chpoint}, a stress
is a \op{mchaml}, and plotting or combining them requires passing the mesh in
the first case and the model in the second.

\subsection*{Operators}

Elementary operators perform a large class of actions:

\begin{itemize}
\item \op{+}, \op{*}, \op{/}, \op{-}, \op{**} are algebraic operators and
  apply equally to numbers and to fields, whether defined at nodes
  (\op{chpoint}) or on elements (\op{mchaml});
\item \op{droi[te]} \dindex{droi[te]}{line}, \op{cerc[le]}
  \dindex{cerc[le]}{circle}, \op{surf[ace]} \dindex{surf[ace]}{surface},
  \op{dall[er]} \dindex{dall[er]}{tiling 2D} construct lines, circles and
  surfaces starting from given points or contours;
\item \op{plus} \sindex{plus} performs a translation by a given vector and
  applies both to a \op{mchaml} and to a geometry;
\item \op{reso[udre]} \dindex{reso[udre]}{solve} solves $KX=Y$, with $K$ a
  symmetric positive definite matrix of \op{rigidite} type and $X$, $Y$ nodal
  fields;
\item \op{vibr[ation]} \dindex{vibr[ation]}{eigenvalue problem, vibration} solves the eigenvalue problem
  $KX-\lambda MX=0$;
\item \op{masq[uer]} \dindex{masq[uer]}{mask, characteristic function} selects
  the values of a field with respect to a user-defined criterion.
\end{itemize}

\noindent Other operators are purely algorithmic:

\begin{itemize}
\item \op{si} \dindex{si}{if}, \op{sinon} \dindex{sinon}{else}, \op{finsi}
  \dindex{finsi}{endif} are the classical control statements, equivalent to
  \texttt{if}, \texttt{else}, \texttt{endif};
\item \op{repe[ter]} \dindex{repe[ter]}{do, start}, \op{fin}
  \dindex{fin}{end} define a loop, corresponding to \texttt{do} \ldots\
  \texttt{enddo};
\item \op{debpr[ocedure]} \dindex{debpr[ocedure]}{subroutine, start} and
  \op{finp[roc]} \dindex{finp[roc]}{subroutine, end} define a
  \emph{procédure}, the equivalent of a subroutine or function.
\end{itemize}

\noindent There are two basic graphical operators:

\begin{itemize}
\item \op{dess[iner]} \dindex{dess[iner]}{draw} plots a curve from lists of
  abscissae and ordinates;
\item \op{trac[er]} \dindex{trac[er]}{plot} displays a mesh, or isolines or
  isocolours of a field defined on a mesh.
\end{itemize}

\noindent A number of operators are complex manipulations of the data
structure and are themselves written in \textsc{gibiane}:

\begin{itemize}
\item \op{pasapas} \dindex{pasapas}{step-by-step, incremental problem}
  incrementally solves nonlinear problems under the small and large strain
  assumptions -- contact, elastoplasticity, and so on;
\item the buckling problem is the eigenvalue problem $KX-\lambda K_{G}X=0$
  with $K$ and $K_{G}$ the stiffness and geometrical stiffness matrices. It is
  solved by the \op{flambage} \sindex{flambage} \emph{procedure}, which
  assembles the matrices and passes the problem to the eigenvalue solver
  (chapter~\ref{chap:eigen}).

  \medskip
  \textbf{Beware of the name.} \op{flambage} must be written out in full.
  The four-letter abbreviation \op{flam} is a \emph{different operator
  altogether} -- it integrates the combustion of a hydrogen--oxygen mixture --
  so \op{flam} in a buckling computation does not fail with a helpful message,
  it calls the wrong thing. This is the one place in \textsc{gibiane} where
  the four-character rule will actively mislead you.
\end{itemize}

\noindent Finally, a feature absent from older editions of these notes:
\texttt{Cast3M} runs in \textbf{parallel}. Setting \op{opti 'PARA' VRAI;} and
\op{opti 'ASSI' n;} (number of assistants) distributes the heavy operations
over up to 128 cores. For the problem sizes in this booklet it makes no
difference, but it matters as soon as you move to three-dimensional models.

In the next chapters we present a series of programming examples illustrating
on the one hand the usage of the \texttt{Cast3M} environment and on the other
hand the programming of classical finite element algorithms. From the
practical point of view, however, do not forget that most classical problem
settings in mechanics already have dedicated operators.

\section*{Summary}
\addcontentsline{toc}{section}{\texorpdfstring{\hskip14mm}{} Summary}

You should now be able to install \texttt{Cast3M}, launch it on a file and
get the hand back at the interactive prompt, and read the notice of any
operator. The three things worth carrying into the next chapters are: the
executable is \texttt{castem26}, not \texttt{castem}; a statement ends at the
\texttt{;} and brackets are compulsory because there is no operator
precedence; and the distinction between a field at the nodes
(\op{chpoint}) and a field on the elements (\op{mchaml}) governs almost
every plotting and combination error you will make.
